\documentclass[12pt,a4paper]{article}

\usepackage[utf8]{inputenc}
\usepackage[portuguese]{babel}
\usepackage[margin=2.5cm]{geometry}
\usepackage{amsmath,amsfonts,amssymb}
\usepackage{graphicx}
\usepackage{xcolor}
\usepackage{hyperref}
\usepackage{listings}
\usepackage{booktabs}
\usepackage{caption}

\definecolor{codegray}{rgb}{0.5,0.5,0.5}
\definecolor{codepurple}{rgb}{0.58,0,0.82}
\definecolor{backcolour}{rgb}{0.05,0.06,0.09}
\definecolor{textcolor}{rgb}{0.9,0.9,0.9}
\definecolor{keywordcolor}{rgb}{0.96,0.62,0.04}
\definecolor{stringcolor}{rgb}{0.2,0.8,0.5}

\lstdefinelanguage{TypeScript}{
  keywords={abstract, any, as, boolean, break, case, catch, class, const, continue, debugger, declare, default, delete, do, else, enum, export, extends, false, finally, for, from, function, get, if, implements, import, in, instanceof, interface, is, keyof, let, module, namespace, new, null, number, object, private, protected, public, readonly, require, return, set, static, string, super, switch, this, throw, true, try, type, typeof, undefined, var, void, while, with, yield},
  morecomment=[l]{//},
  morecomment=[s]{/*}{*/},
  morestring=[b]',
  morestring=[b]"
}

\lstset{
  backgroundcolor=\color{backcolour},
  commentstyle=\color{codegray},
  keywordstyle=\color{keywordcolor}\bfseries,
  numberstyle=\tiny\color{codegray},
  stringstyle=\color{stringcolor},
  basicstyle=\ttfamily\footnotesize\color{textcolor},
  breakatwhitespace=false,
  breaklines=true,
  captionpos=b,
  keepspaces=true,
  numbers=left,
  numbersep=5pt,
  showspaces=false,
  showstringspaces=false,
  showtabs=false,
  tabsize=2
}

\hypersetup{
    colorlinks=true,
    linkcolor=blue,
    filecolor=magenta,      
    urlcolor=cyan,
    pdftitle={Relatorio de Progresso e Arquitetura - VTT RPG},
}

\title{\textbf{Relatório Técnico de Progresso e Arquitetura}\\ Sistema VTT RPG Web de Alta Fidelidade (D\&D 5e)}
\author{\textbf{Equipe de Desenvolvimento VTT RPG}}
\date{\today}

\begin{document}

\maketitle

\begin{abstract}
Este relatório apresenta o detalhamento técnico do progresso e da arquitetura do software para a plataforma web de VTT (\textit{Virtual Tabletop}) inspirada nas mecânicas do D\&D 5e e na interface visual do jogo \textit{Baldur's Gate 3}. São documentados os módulos de criação de personagens com mecânica de \textit{Point Buy}, mapa tático baseado em grelha com cálculo de alcance por distância de Chebyshev, motor de interpretação de rolagens de dados (\textit{parser} regex com suporte a acertos/falhas críticas), log de batalha em tempo real sincronizado via WebSockets e o painel de controle do Mestre. O documento aborda a modelagem relacional do banco de dados no Supabase (PostgreSQL) e as políticas de segurança \textit{Row Level Security} (RLS) imutáveis para integridade de dados.
\end{abstract}

\tableofcontents
\newpage

\section{Introdução e Escopo do Projeto}
O objetivo deste projeto é o desenvolvimento de uma plataforma web de \textit{Virtual Tabletop} (VTT) focada em proporcionar uma experiência imersiva, moderna e fluida para partidas de RPG de Mesa baseadas na 5ª Edição do \textit{Dungeons \& Dragons} (D\&D 5e).

\subsection{Stack Tecnológica Escolhida}
\begin{itemize}
    \item \textbf{Frontend Framework:} Next.js (React 18 / App Router) com TypeScript.
    \item \textbf{Estilização e UI:} Tailwind CSS.
    \item \textbf{Backend \& Persistência:} Supabase (PostgreSQL).
    \item \textbf{Sincronização Multiplayer:} Supabase Realtime (WebSockets).
    \item \textbf{Segurança:} Políticas de \textit{Row Level Security} (RLS).
\end{itemize}

\section{Arquitetura da Base de Dados e Segurança (RLS)}
A persistência de dados foi projetada sobre um modelo relacional no PostgreSQL.

\subsection{Modelo Relacional}
\begin{enumerate}
    \item \texttt{campaigns}: Armazena os dados da sessão.
    \item \texttt{tokens}: Representa os elementos do mapa tático.
    \item \texttt{room\_events}: Log unificado do chat e rolagens.
    \item \texttt{room\_members}: Associação N:N entre utilizadores e campanhas.
\end{enumerate}

\section{Estado Atual do Projeto e Próximos Passos}
\begin{table}[h!]
\centering
\caption{Status de Desenvolvimento dos Módulos do VTT}
\begin{tabular}{lll}
\toprule
\textbf{Módulo / Funcionalidade} & \textbf{Status} & \textbf{Tecnologias Envolvidas} \\
\midrule
Criação de Personagens (Point Buy) & Concluído & React, TypeScript, Tailwind CSS \\
HUD Ativo do Jogador & Concluído & React, State Management \\
Mapa Tático 2D \& Grid & Concluído & React, Distância de Chebyshev \\
Base de Dados Relacional \& RLS & Concluído & Supabase, PostgreSQL, RLS \\
Motor de Dados (Dice Parser) & Concluído & TypeScript, Regex Engine \\
Sincronização Multiplayer Realtime & Concluído & Supabase Realtime, WebSockets \\
Painel do Mestre (DM Dashboard) & Concluído & React, Supabase API \\
Integração da Tela Única (GameRoom)& Concluído & Next.js App Router, Tailwind \\
Gerador de Avatares por IA & Pendente & OpenAI DALL-E 3 / Replicate API \\
\bottomrule
\end{tabular}
\end{table}

\begin{thebibliography}{9}
\bibitem{srd51} Wizards of the Coast. \textit{System Reference Document 5.1 (SRD5)}. 2016.
\bibitem{supabaseDocs} Supabase Inc. \textit{Supabase Documentation}. 2026.
\bibitem{nextjsDocs} Vercel. \textit{Next.js App Router Documentation}. 2026.
\end{thebibliography}

\end{document}
